\paragraph{Passage to the original limiting group.}
Keep the normalized point, odd family, auxiliary tower and its original
limiting quotient. Put $A=k[[t]]$, $A_m=A/(t^m)$ and
$D=n(d+q_{\rm aux})$. All coefficient modules, ordinary witnesses,
diagonal restrictions, inertia conditions and polarizations below are
the original ones.

\begin{proposition}[The limiting Selmer conditions]
\label{prop:selmer-limit-comparison}
The finite-level estimates of Propositions
\ref{prop:selmer-first-global-bound} and
\ref{prop:selmer-second-global-bound} give the following conditions on
the original limiting group.
\begin{enumerate}
\item At most $D$ integral first-order Selmer classes are independent
modulo the prescribed polarized boundaries.
\item Every $g+1$ tuples $p_{ij}$ of integral even Selmer cocycles admit
an integral second-order Selmer cocycle $\Psi$ and coefficients
$a_i=t^{C_2}c_i$, with some $c_i$ a unit, such that
$\Psi_{\epsilon'_j}=\sum_i a_i p_{ij}$.
\item Every integral odd Selmer cocycle is in the span of the original
$\kappa_j$ modulo the allowed boundaries after multiplication by a
power of $t$.
\end{enumerate}
The corresponding Laurent statements hold for the bounded continuous
cocycles in the local deformation presentation.
\end{proposition}

\begin{proof}
Write $P$ for the original common profinite cover, with quotient maps
$\pi_N,\pi_\infty$. Its limit data assert, for every open normal
$V\subset P$, that
$(\ker\pi_N)V=(\ker\pi_\infty)V$ for almost every $N$ in the
chosen free ultrafilter. The original base and local maps agree modulo
$V$ in the same sense. The open subgroup $P^+$ is topologically
finitely generated, so its continuous cocycles with any fixed finite
coefficient module form a finite set: their generator values determine
them. Each such cocycle is constant modulo some $V$. Intersecting the
finitely many good sets therefore identifies the finite Selmer sets on
the two quotients simultaneously. Their local restrictions agree and
the same witness matrices work. The boundary sets also agree because
the coefficient actions agree. Thus the finite Selmer quotients are
identified, preserving both constants and every local condition.

For the first condition, let $z_1,\ldots,z_e$ be independent integral
classes. Evaluate $\sum_i c_i z_i-dY$ on fixed topological generators,
with $Y$ ranging over the finite free module of allowed skew primitives.
The kernel of this finite linear map projects to zero in $A^e$.
Smith normal form gives $b\geq0$ such that a boundary congruence
modulo $t^m$ forces every $c_i\in t^{\max(m-b,0)}A$.
Allowed skew primitives modulo $t^m$ lift integrally by lifting one
entry in each skew pair; odd characteristic makes the diagonal zero.
Consequently
\[
 |k|^{e(m-b)}\leq
 \#(\operatorname{Sel}^1_{\infty,m}/B^1_{\infty,m})
 \leq |k|^{mD+C_1}\qquad(m\geq b).
\]
Letting $m$ increase gives $e\leq D$, with the finite congruence
defect $b$ retained.

For the second condition, at each $m\geq1$ keep all solutions
consisting of a primitive coefficient vector in $A_m^{g+1}$,
a second-order cocycle and its ordinary witness matrices. These are
nonempty finite sets, and reduction preserves them. Their inverse
limit is nonempty even when the reduction maps are not surjective:
every finite compatibility system comes from a sufficiently high
level. The resulting cocycle is continuous because each finite
reduction is continuous. Separatedness of $A$ preserves every
cocycle, polarization, inertia and ordinary equation. A coefficient
nonzero modulo $t$ remains a unit; hence the vector $t^{C_2}c$ is
nonzero. This proves the assertion with the original loss $C_2$.

The third condition is the resultant descent following
Lemma~\ref{lem:selmer-block-resultant}. The auxiliary inertia elements
normally generate the limiting kernel, and their odd cocycle values
commute with separated Frobenius blocks, so they vanish. Polarization
also gives invariance under complex conjugation. The cocycle descends
to the original base group and its chosen odd generators.

Finally, a common power of $t$ makes any finite family of bounded
Laurent cocycles and their finitely many local witness matrices
integral. Apply the three integral assertions and invert that power.
The full coefficient action, all local equations and the nonzero
coefficient condition are preserved \cite{alpo-selmer-limit-comparison}.
\end{proof}

This supplies the three cohomological inputs of the local presentation
and incidence argument. The actual arithmetic models and the patched
automorphic support of dimension at least $1+D$ remain to be supplied.
The rank conclusion does not bound integral torsion uniformly over
moving points.
